% Part: lambda-calculus
% Chapter: introduction
% Section: term-revisited

\documentclass[../../../include/open-logic-section]{subfiles}

\begin{document}

\olfileid[hi]{lam}{syn}{tr}
\olsection{$\alpha$-तुल्यता वर्गों के रूप में पद}

अब से हम पदों पर $\alpha$-तुल्यता तक विचार करेंगे। इसका अर्थ है कि जब हम
कोई पद लिखते हैं, तब हमारा आशय उस $\alpha$-तुल्यता वर्ग से होता है जिसमें
वह पद है। उदाहरणार्थ, हम $\lambd[a][\lambd[b][a c]]$ को उसके
$\alpha$-तुल्य सभी पदों के समुच्चय के लिए लिखते हैं, जैसे
$\lambd[a][\lambd[b][a c]]$, $\lambd[b][\lambd[a][b c]]$, आदि।

इसके अतिरिक्त, पिछले खंडों में $N, Q$ जैसे अक्षर किसी पद को दर्शाते थे,
पर अब से हम उनसे किसी वर्ग को दर्शाएँगे; आगे पदों के बजाय यही वर्ग हमारे
अध्ययन की विषय-वस्तु होंगे। $x, y$ जैसे अक्षर अब भी चर को दर्शाते हैं।

हम वर्ग~$M$ के किसी स्वेच्छ !!{element} को दर्शाने के लिए संकेत-लिपि
$\rep{M}$ भी अपनाते हैं; और एक से अधिक की आवश्यकता होने पर $\rep{M}[0]$,
$\rep{M}[1]$, आदि लिखते हैं।

अपनी भाषा सरल रखने के लिए हम पदों वाली संकेत-लिपि फिर उपयोग करेंगे। वर्गों
पर हमारी निम्न परिभाषा है:
\begin{defn}
  \begin{enumerate}
  \item $\lambd[x][N]$ उस वर्ग के रूप में परिभाषित है जिसमें
    $\lambd[x][\rep{N}]$ है।
  \item $PQ$ उस वर्ग के रूप में परिभाषित है जिसमें $\rep{P}\rep{Q}$ है।
  \end{enumerate}
\end{defn}

यह देखना कठिन नहीं कि ये सुपरिभाषित हैं, क्योंकि $\alpha$-परिवर्तन संगत
है।

\begin{defn} \ollabel{def:fv}
  किसी $\alpha$-तुल्यता वर्ग~$M$ के \emph{मुक्त चर}, अथवा $FV(M)$, को
  $FV(\rep{M})$ परिभाषित करते हैं।
\end{defn}

यह सुपरिभाषित है, क्योंकि \olref[alp]{thm:fv} में दिखाए अनुसार
$FV(\rep{M}[0]) = FV(\rep{M}[1])$।

वर्गों में प्रतिस्थापन के लिए भी हम वही संकेत-लिपि पुनः उपयोग करते हैं:
\begin{defn} \ollabel{def:sub}
  ~$M$ में $y$ के स्थान पर $R$ का \emph{प्रतिस्थापन}, अथवा
  $\Subst{M}{R}{y}$, ऐसा कोई भी $\rep{M}$ और $\rep{R}$ लेकर
  $\Subst{\rep{M}}{\rep{R}}{y}$ परिभाषित है जिसके लिए यह प्रतिस्थापन
  परिभाषित हो।
\end{defn}

\olref[alp]{cor:sub} में दिखाए अनुसार यह भी सुपरिभाषित है।

ध्यान दें कि यह परिभाषा हमारे तर्क को कितना सरल करती है। उदाहरणार्थ:
\begin{align}
  \Subst{\lambd[x][x]}{y}{x} & =\ollabel{eq:1}\\
  &= \Subst{\lambd[z][z]}{y}{x} \ollabel{eq:2}\\
  &= \lambd[z][\Subst{z}{y}{x}] \\
  &= \lambd[z][z]
\end{align}

यदि हम अब भी इसे पदों पर प्रतिस्थापन मानें, तो \olref{eq:1} अपरिभाषित है;
किंतु पहले बताए अनुसार अब हम इसे वर्गों पर प्रतिस्थापन मानते हैं। इसीलिए
\olref{eq:2} संभव है: हम $\lambd[x][x]$ को $\lambd[z][z]$ से बदल सकते हैं,
क्योंकि वे एक ही वर्ग में हैं।

इसी कारण अब से हम मानेंगे कि हमारे चुने प्रतिनिधि प्रतिस्थापन के लिए
आवश्यक शर्तों को सदा संतुष्ट करते हैं। उदाहरणार्थ, जब हमें
$\Subst{\lambd[x][N]}{R}{y}$ दिखेगा, तब हम मानेंगे कि प्रतिनिधि
$\lambd[x][N]$ ऐसा चुना गया है कि $x \neq y$ और $x \notin FV(R)$।

चूँकि $\lambd[x][x]$ को ``वर्ग'' कहना कुछ विचित्र है, इसलिए अब से इन्हें
$\Lambda$-पद (अथवा इस भाग के शेष अंश में केवल ``पद'') कहें, ताकि इन्हें
हमारे परिचित $\lambda$-पदों से अलग किया जा सके।

\begin{editorial}
  अभी हम पदों को पूरी तरह विदा नहीं कर सकते: $\Lambda$-पदों की पूरी
  परिभाषा $\lambda$-पदों पर आधारित है, और हमने $\Lambda$-पदों पर फलन
  परिभाषित करने की कोई विधि नहीं दी है। इसका अर्थ है कि ऐसे सभी फलन पहले
  $\lambda$-पदों पर परिभाषित करने और फिर $\Lambda$-पदों पर ``प्रक्षेपित''
  करने होंगे, जैसा हमने प्रतिस्थापन के लिए किया था। फिर भी हम मानते हैं कि
  पाठक सहज रूप से समझ सकता है कि $\Lambda$-पदों पर फलन कैसे परिभाषित किए
  जा सकते हैं।
\end{editorial}
\end{document}
